\sgaStableTarget{sga6.e05.def.3-9}\textbf{Definition 3.9.} \uline{Let} @@M0@@ \uline{be a} @@M1@@-\uline{ring}; \uline{by an} augmented @@M2@@-@@M3@@-algebra \uline{we mean the data of a} @@M4@@-@@M5@@-\uline{algebra} @@M6@@ \uline{and a} @@M7@@-\uline{homomorphism of} @@M8@@-\uline{algebras} @@M9@@; a homomorphism of augmented @@M10@@-@@M11@@-algebras \uline{is a} @@M12@@-\uline{homomorphism compatible with the augmentations}.

\sgaStableTarget{sga6.e05.ex.3-9-1}\textbf{Example 3.9.1.}: Let @@M13@@ be a commutatively locally ringed topos, and @@M14@@ the Grothendieck group of the category of sheaves of locally free @@M15@@-modules of finite type on @@M16@@. We have seen that @@M17@@ can be endowed with a pre-@@M18@@-ring structure (2.2.b, since every direct summand of a free module of finite type is then locally free of finite type) and that @@M19@@ is canonically endowed with a binomial ring structure. We define a @@M20@@-homomorphism from @@M21@@ to @@M22@@ as follows: let @@M23@@ be a section of the constant sheaf @@M24@@ on @@M25@@; it defines a partition of the final object of @@M26@@ into a sum of disjoint opens @@M27@@, the section @@M28@@ taking the value @@M29@@ on the open @@M30@@; let @@M31@@ be the sheaf equal to @@M32@@ on @@M33@@ for @@M34@@, and to @@M35@@ on @@M36@@ for @@M37@@, and let @@M38@@ be the sheaf equal to @@M39@@ on @@M40@@ for @@M41@@, and to @@M42@@ on @@M43@@ for @@M44@@; we associate to @@M45@@ the element @@M46@@ of @@M47@@.

It is immediate to verify that this does give a @@M48@@-homomorphism from @@M49@@ to @@M50@@; we can then define the augmentation from @@M51@@ to @@M52@@ by assigning to every locally free sheaf of finite type its rank, which is a section of @@M53@@; additivity of rank implies that this defines a function on @@M54@@ with values in @@M55@@. We have thus endowed @@M56@@ with an augmented pre-@@M57@@-algebra structure. This is always the structure referred to below.

